% Main table: `plain-bold` (design-axes.md axis 8; P2, P5).
% Plain booktabs, no tints, no vertical rules. \best marks the winning cell
% in each metric column; the direction is stated once per column header.
% A Group/Dataset multirow may group rows. Caption ABOVE.
%
% Reporting uncertainty (references/table-grammar.md, "Reporting
% uncertainty" -- read that section before filling this in): every metric
% cell, for every arm, carries `mean [lo, hi]`, lo/hi = min/max of that arm's
% own per_split[condition]['values'] -- bracket, never `mean $\pm$ std`,
% because the range is observed and generally asymmetric, not a symmetric
% half-width. Both real draws of this grammar (S16, P247) fit the full
% per-cell treatment at 8pt with \acfit added defensively (the unbracketed
% version did not always need it; the bracketed one can cross \linewidth
% even for a "plain" 4-5-column table). If this grammar is drawn with an
% Overall/Mean column (P247's shape does; the block below does not), that
% column stays a point estimate: its own `std`, if the aggregate reports
% one, is spread ACROSS conditions, not repeated-run noise.
\begin{acwidetable}
  \caption{\textbf{\slot{Main comparison, named}.} \slot{one sentence on the
    protocol, the estimator, and the precision reported}; bracketed values
    are $[\min,\max]$ over the repeated runs behind each mean.}
  \label{tab:main-results}
  \vspace{-0.2cm}
  \centering
  {\fontsize{8pt}{10pt}\selectfont
   % \acfit, not \resizebox: shrink to the float width when the table is too
   % wide, and leave it alone when it is not (references/table-grammar.md).
   % A plain-looking table without it was the one that most often surprised a
   % writer by crossing \linewidth once every cell grew an interval.
   \acfit{%
   \begin{tabular}{l l c c c c}
     \toprule
     \textbf{Group} & \textbf{Arm}
       & \textbf{\slot{metric 1}\down} & \textbf{\slot{metric 2}\down}
       & \textbf{\slot{metric 3}\up} & \textbf{\slot{metric 4}\up} \\
     \midrule
     \multirow{3}{*}{\textbf{\slot{Group A}}}
       & \slot{baseline one} & \slot{v} \slot{lo, hi} & \slot{v} \slot{lo, hi}
                             & \slot{v} \slot{lo, hi} & \slot{v} \slot{lo, hi} \\
       & \slot{baseline two} & \slot{v} \slot{lo, hi} & \slot{v} \slot{lo, hi}
                             & \slot{v} \slot{lo, hi} & \slot{v} \slot{lo, hi} \\
       & \slot{primary arm}  & \best{\slot{v} \slot{lo, hi}} & \best{\slot{v} \slot{lo, hi}}
                             & \best{\slot{v} \slot{lo, hi}} & \best{\slot{v} \slot{lo, hi}} \\
     \midrule
     \multirow{3}{*}{\textbf{\slot{Group B}}}
       & \slot{baseline one} & \slot{v} \slot{lo, hi} & \slot{v} \slot{lo, hi}
                             & \slot{v} \slot{lo, hi} & \slot{v} \slot{lo, hi} \\
       & \slot{baseline two} & \slot{v} \slot{lo, hi} & \slot{v} \slot{lo, hi}
                             & \slot{v} \slot{lo, hi} & \slot{v} \slot{lo, hi} \\
       & \slot{primary arm}  & \best{\slot{v} \slot{lo, hi}} & \best{\slot{v} \slot{lo, hi}}
                             & \best{\slot{v} \slot{lo, hi}} & \best{\slot{v} \slot{lo, hi}} \\
     \bottomrule
   \end{tabular}}}
\end{acwidetable}
